%% ma: L11-B15-0010
%% lop: 11
%% chuong: 5
%% bai: 15
%% dang: 11.15.5
%% loai: TL
%% muc: VD
%% dap_an: "-"
%% nguon: "tai-lieu-hoc-tap-mon-toan-lop-11-theo-chuong-trinh-giao-duc-pho-thong-2018.pdf (Lớp Toán Thầy Duy, Phạm Lê Duy, THPT Chu Văn An), Chương 3 Bài 1, Dạng 4, Ví dụ 10 (tr.313)"
%% kien_thuc: "Lớp 11 Bài 15 (giới hạn của dãy số)"
%% trang_thai: cho-duyet
%% ly_do: "Dạng mới đề xuất (11.15.x), chờ giáo viên duyệt dạng; tài liệu gốc không có lời giải mẫu, đáp án do kho tự giải và kiểm bằng sympy."
%% ghi_chu: "Ý b) cần nhân lượng liên hợp bậc ba (hằng đẳng thức hiệu hai lập phương), SGK Bài 15 chưa nêu tường minh nhưng là phép biến đổi đại số cơ bản."
\begin{ex}
Tính các giới hạn sau:

a) $\lim\dfrac{\sqrt{4n^2+1}+2n-1}{\sqrt{n^2+4n+1}+n}$;

b) $\lim\dfrac{n^2+\sqrt[3]{1-n^6}}{\sqrt{n^4+1}+n^2}$.
\loigiai{\sloppy a) Chia tử và mẫu cho $n$: \[\dfrac{\sqrt{4+\frac1{n^2}}+2-\frac1n}{\sqrt{1+\frac4n+\frac1{n^2}}+1}\to\dfrac{2+2}{1+1}=2.\]

b) Đặt $\sqrt[3]{1-n^6}=-n^2\sqrt[3]{1-\frac1{n^6}}$. Tử \[=n^2\left(1-\sqrt[3]{1-\frac1{n^6}}\right)=n^2\cdot\dfrac{\frac1{n^6}}{1+\sqrt[3]{1-\frac1{n^6}}+\sqrt[3]{\left(1-\frac1{n^6}\right)^2}}\] (dùng $1-a^3=(1-a)(1+a+a^2)$ với $a=\sqrt[3]{1-\frac1{n^6}}$), nên tử $\sim\dfrac{1}{3n^4}\to0$. Mẫu $=n^2\left(\sqrt{1+\frac1{n^4}}+1\right)\sim2n^2$. Vậy giới hạn bằng $0$.}
\end{ex}
